\subsection{CE2 normal-form compression versus a blind derivation}
\label{subsec:ce2-blindness}

A source-level audit distinguishes two notions that had previously been conflated.  A
\emph{runtime-blind} predictor does not consult the target answer table while producing its
predictions; a \emph{provenance-blind} predictor additionally uses no coefficient or rule
that was fitted or selected from those target answers.  A genuine no-answer-table replay
requires both.

The current CE2 normal-form machinery does not yet meet that standard.  In particular,
\verb|_derive_tables_via_normal_form()| explicitly reads
\verb|_simple_family_sign_map()| when constructing its constant-line tables, while
\verb|_derive_naive_tables()| uses the same map to supply the domain keys.  More decisively,
the source documents the delta-polynomial coefficients as fits obtained from the full
$864$-entry simple-family sign dataset.  Hence the current result is accurately described as
\[
 \boxed{\text{structural compression verified; independent blind derivation open}.}
\]
This correction does not invalidate the native $27\times3$ CE2 sparse rows, the
Heisenberg/metaplectic support laws, or the empirical correctness of the compressed
predictor.  It changes only the evidentiary tier of the global closed form.  The missing
certificate must derive the seed and delta laws from the bracket/cochain structure,
enumerate the domain without answer keys, freeze all predictions, and only then compare
against the committed sparse-solution table.
